For $u, u' \in U$, $u$ and $u'$ anticommute with $e$ (\claimref{std-algebra.orthogonal-anticommute}), so $(ue)(u'e) = -u u' e^{2} = -s\,u u'$ and $(ue)(u'e) + (u'e)(ue) = -s(uu' + u'u) = -2s\,B(u, u')$. Hence $(ue)^{2} = -s\,Q(u)$, and the universal property (\claimref{std-algebra.existence-uniqueness}) gives a homomorphism $\phi : \Cl(U, -sQ) \to \Cl(V, Q)$ with image in $\Cl^{+}$ (\claimref{std-algebra.z2-grading}). Take an orthogonal basis $u_1, \dots, u_{n-1}$ of $U$. For $I = \{i_1 < \dots < i_k\}$, $\phi(u_{i_1}) \cdots \phi(u_{i_k}) = \pm s^{\lfloor k/2 \rfloor} u_I e^{k \bmod 2}$, so the image contains every even blade $u_I$ ($k$ even) and $u_I e$ ($k$ odd) of the orthogonal basis $u_1, \dots, u_{n-1}, e$, which span $\Cl^{+}$. Both sides have dimension $2^{n-1}$ (\claimref{std-algebra.blade-basis}), so $\phi$ is an isomorphism. With $Q(e) = -1$ the signature of $(U, Q|_U)$ is $(p, q - 1, r)$ and is unchanged; with $Q(e) = 1$ it is $(p - 1, q, r)$ and the sign flip makes it $(q, p - 1, r)$.
